\documentclass[border=0pt]{standalone}
\usepackage{tikz}
\usetikzlibrary{knots,hobby,intersections,calc}
\definecolor{avebg}{HTML}{0B0E12}
\definecolor{cyanp}{HTML}{34E1EE}
\definecolor{fluxp}{HTML}{FF3A66}
\definecolor{fluxhi}{HTML}{FFA6BC}
\def\TRE{plot[domain=0:360,samples=160,smooth,variable=\t]({0.80*(sin(\t)+2*sin(2*\t))},{0.80*(cos(\t)-2*cos(2*\t))})}
\begin{document}
\begin{tikzpicture}
\clip (-5.35,-5.35) rectangle (5.35,5.35);
\fill[avebg] (-5.7,-5.7) rectangle (5.7,5.7);

% ===== anisotropic lensed K4 lattice (4-fold, NOT radially symmetric) =====
\foreach \i in {0,...,22}{\foreach \j in {0,...,22}{
  \pgfmathsetmacro{\gx}{(\i-11)*0.52}\pgfmathsetmacro{\gy}{(\j-11)*0.52}
  \pgfmathsetmacro{\r}{sqrt(\gx*\gx+\gy*\gy)+0.001}
  \pgfmathsetmacro{\th}{atan2(\gy,\gx)}
  \pgfmathsetmacro{\f}{max(0.34, 1 - 1.7*(1+0.45*cos(4*\th))/(\r*\r+1.5))}
  \coordinate (n-\i-\j) at (\gx*\f,\gy*\f);}}
% bonds (brighter toward the lensing centre)
\foreach \i in {0,...,21}{\foreach \j in {0,...,22}{
  \pgfmathsetmacro{\gx}{(\i-10.5)*0.52}\pgfmathsetmacro{\gy}{(\j-11)*0.52}\pgfmathsetmacro{\r}{sqrt(\gx*\gx+\gy*\gy)}
  \pgfmathsetmacro{\op}{min(0.55,0.08+0.6/(\r+1.0))}
  \pgfmathtruncatemacro{\ii}{\i+1}\draw[cyanp,opacity=\op,line width=0.45pt](n-\i-\j)--(n-\ii-\j);}}
\foreach \i in {0,...,22}{\foreach \j in {0,...,21}{
  \pgfmathsetmacro{\gx}{(\i-11)*0.52}\pgfmathsetmacro{\gy}{(\j-10.5)*0.52}\pgfmathsetmacro{\r}{sqrt(\gx*\gx+\gy*\gy)}
  \pgfmathsetmacro{\op}{min(0.55,0.08+0.6/(\r+1.0))}
  \pgfmathtruncatemacro{\jj}{\j+1}\draw[cyanp,opacity=\op,line width=0.45pt](n-\i-\j)--(n-\i-\jj);}}
% nodes
\foreach \i in {0,...,22}{\foreach \j in {0,...,22}{
  \pgfmathsetmacro{\gx}{(\i-11)*0.52}\pgfmathsetmacro{\gy}{(\j-11)*0.52}\pgfmathsetmacro{\r}{sqrt(\gx*\gx+\gy*\gy)}
  \pgfmathsetmacro{\op}{min(0.9,0.1+0.75/(\r+0.8))}
  \fill[cyanp,opacity=\op](n-\i-\j) circle(0.032);}}

% ===== glassy torus annulus (the lensing mass), top-down =====
\fill[even odd rule,cyanp,opacity=0.18] (0,0) circle(2.45) (0,0) circle(1.05);
\draw[cyanp,opacity=0.40,line width=3pt] (0,0) circle(1.75);     % tube top sheen
\draw[cyanp,opacity=0.75,line width=1.1pt] (0,0) circle(2.45);   % outer edge
\draw[cyanp,opacity=0.75,line width=1.1pt] (0,0) circle(1.05);   % inner edge
\draw[white,opacity=0.22,line width=2pt] (118:1.75) arc(118:165:1.75); % glass highlight

% ===== trefoil mark — stylized real-space rendering of the (2,3) =====
% NB: canonically the (2,3) is a PHASE-SPACE knot in the bond-LC (V_inc,V_ref)
%     phasor (theory.md:16); the real-space electron is the 0_1 unknot envelope.
%     Drawn here as a real over/under knot for visual legibility only — it is
%     NOT a claim of real-space flux-tube topology.
\draw[fluxp,opacity=0.06,line width=20pt,line cap=round] \TRE;   % outer bloom
\draw[fluxp,opacity=0.13,line width=12pt,line cap=round] \TRE;   % glow halo
\begin{knot}[consider self intersections=true,clip width=6,end tolerance=1pt,
  flip crossing=2]
  \strand[fluxp,line width=5pt,line cap=round] \TRE;
\end{knot}
\begin{knot}[consider self intersections=true,clip width=6,end tolerance=1pt,
  flip crossing=2]
  \strand[fluxhi,line width=1.6pt,line cap=round] \TRE;        % highlight core
\end{knot}
\end{tikzpicture}
\end{document}
